\section{Conclusion}
We examine how language models respond to repeated user challenges and whether baseline uncertainty explains their susceptibility to sycophancy. Across GPT, Claude, and Gemini models on MMLU-Pro, GPQA-Diamond, and HLE, uncertain questions generally exhibit higher flip rates, although the association varies across models and benchmarks and does not survive multiple-comparison correction. We do not find consistent evidence that uncertainty amplifies the effect of conversational pressure. Resistance and capitulation also vary across model--benchmark pairs rather than reflecting fixed model traits. Further, correctness can decline sharply while confidence inferred from reasoning traces remains high or increases. These findings suggest that baseline uncertainty alone is insufficient to characterize susceptibility to sycophancy. Evaluating multi-turn response trajectories alongside correctness and estimated confidence provides a more complete account of model reliability under adversarial conversational pressure.
